\documentclass[11pt,a4paper]{article}

% ── packages ─────────────────────────────────────────────────────────────────
\usepackage[margin=2.5cm]{geometry}
\usepackage{amsmath}
\usepackage{amssymb}
\usepackage{booktabs}
\usepackage{hyperref}
\usepackage{xcolor}
\usepackage{longtable}
\usepackage{array}
\usepackage{graphicx}
\usepackage{fancyhdr}
\usepackage{titlesec}
\usepackage{multirow}
\usepackage{enumitem}

% ── colours ──────────────────────────────────────────────────────────────────
\definecolor{apexblue}{RGB}{0,51,102}
\definecolor{apexgray}{RGB}{100,100,100}
\definecolor{apexred}{RGB}{180,30,30}

% ── hyperref setup ───────────────────────────────────────────────────────────
\hypersetup{
  colorlinks=true,
  linkcolor=apexblue,
  citecolor=apexblue,
  urlcolor=apexblue,
  pdftitle={MDD: PFE / CCR Exposure Engine -- APEX-MDL-0014},
  pdfauthor={Apex Global Bank, Quantitative Research Division}
}

% ── header / footer ──────────────────────────────────────────────────────────
\pagestyle{fancy}
\fancyhf{}
\lhead{\textcolor{apexblue}{\small\textbf{APEX-MDL-0014} \textbar\ PFE / CCR Exposure Engine v1.0}}
\rhead{\textcolor{apexgray}{\small RESTRICTED --- Model Risk Sensitive}}
\lfoot{\textcolor{apexgray}{\small Apex Global Bank \textbar\ Quantitative Research Division}}
\rfoot{\textcolor{apexgray}{\small Page \thepage}}
\renewcommand{\headrulewidth}{0.4pt}
\renewcommand{\footrulewidth}{0.4pt}

% ── section style ────────────────────────────────────────────────────────────
\titleformat{\section}{\large\bfseries\color{apexblue}}{{\thesection.}}{0.5em}{}[\titlerule]
\titleformat{\subsection}{\normalsize\bfseries\color{apexblue}}{{\thesubsection}}{0.5em}{}

% ── misc ─────────────────────────────────────────────────────────────────────
\setlength{\parskip}{0.5em}
\setlength{\parindent}{0pt}

% ─────────────────────────────────────────────────────────────────────────────
\begin{document}
% ─────────────────────────────────────────────────────────────────────────────

% ── title block ──────────────────────────────────────────────────────────────
\begin{center}
  {\LARGE\bfseries\color{apexblue} Model Development Document}\\[0.4em]
  {\Large\bfseries PFE / CCR Exposure Engine}\\[0.3em]
  {\large Version 1.0}\\[0.6em]
  {\normalsize Apex Global Bank \textbar\ Quantitative Research Division}\\[0.2em]
  {\small\color{apexgray} SR 11-7 Section Reference: Model Development \textbar\ March 2026}
\end{center}

\vspace{1em}
\textcolor{apexred}{\textbf{Classification: RESTRICTED --- Model Risk Sensitive}}
\vspace{1em}

\hrule
\vspace{1em}

% ─────────────────────────────────────────────────────────────────────────────
\section{Document Metadata}
% ─────────────────────────────────────────────────────────────────────────────

\begin{table}[h!]
\centering
\begin{tabular}{@{}ll@{}}
\toprule
\textbf{Field} & \textbf{Value} \\
\midrule
Model Name           & PFE / CCR Exposure Engine \\
Model ID             & APEX-MDL-0014 \\
Version              & 1.0 \\
Status               & In Validation \\
Risk Tier            & Tier 1 \\
Regulatory Framework & Basel III SA-CCR; IFRS 13 Fair Value; SR 11-7 \\
Date of Development  & March 2026 \\
Business Owner       & Head of Global Markets Trading (James Okafor) \\
Model Developer      & Quantitative Research --- XVA Team \\
MRO Review           & Pending --- Dr.\ Rebecca Chen \\
CRO Approval         & Pending --- Dr.\ Priya Nair \\
\bottomrule
\end{tabular}
\caption{Document metadata}
\end{table}

\subsection{Version History}

\begin{table}[h!]
\centering
\begin{tabular}{@{}lll@{}}
\toprule
\textbf{Version} & \textbf{Date} & \textbf{Change Summary} \\
\midrule
0.1 & Jan 2026 & Initial netting framework; EPE/ENE profiles only \\
0.8 & Feb 2026 & Collateral offset with MPoR; SA-CCR add-on module \\
1.0 & Mar 2026 & Production candidate; submitted for validation \\
\bottomrule
\end{tabular}
\caption{Version history}
\end{table}

% ─────────────────────────────────────────────────────────────────────────────
\section{Executive Summary}
% ─────────────────────────────────────────────────────────────────────────────

\subsection{Model Purpose}

The PFE / CCR Exposure Engine is the foundational simulation model for all
counterparty credit risk (CCR) measurement at Apex Global Bank. It generates
the exposure profiles --- Expected Positive Exposure (EPE), Expected Negative
Exposure (ENE), and Potential Future Exposure (PFE) --- that are consumed by
three downstream Tier~1 models:

\begin{itemize}[noitemsep]
  \item \textbf{CVA} (APEX-MDL-0010): uses EPE to compute the market value of
        counterparty default risk;
  \item \textbf{FVA} (APEX-MDL-0011): uses EPE/ENE to compute the funding cost
        of uncollateralised positions;
  \item \textbf{MVA} (APEX-MDL-0012): uses simulated future Initial Margin
        profiles derived from the same Monte Carlo paths.
\end{itemize}

In addition, the model produces:

\begin{itemize}[noitemsep]
  \item \textbf{PFE for credit limit management}: the 97.5th-percentile
        exposure used by the Credit Risk function to set and monitor
        counterparty credit limits;
  \item \textbf{SA-CCR capital}: the Basel~III standardised-approach CCR
        exposure-at-default (EAD) used for regulatory capital calculation.
\end{itemize}

Because this model is upstream of multiple Tier~1 models and directly feeds
regulatory capital reporting, it is classified as \textbf{Tier~1}.

\subsection{Output Description}

\begin{table}[h!]
\centering
\begin{tabular}{@{}lll@{}}
\toprule
\textbf{Output} & \textbf{Consumer} & \textbf{Frequency} \\
\midrule
EPE profile (time series) & CVA, FVA, MVA models       & Daily \\
ENE profile (time series) & FVA model                   & Daily \\
PFE at 97.5\% (time series) & Credit limit system       & Daily \\
Effective EPE (Basel)       & Regulatory capital (IMM)  & Daily \\
SA-CCR EAD by netting set   & Capital reporting, RWA    & Daily \\
Netting benefit report      & Credit / legal            & Monthly \\
\bottomrule
\end{tabular}
\caption{Model outputs and consumers}
\end{table}

\subsection{Key Assumptions}

\begin{enumerate}[noitemsep]
  \item Market factors evolve under the risk-neutral measure (for XVA) and the
        real-world measure (for PFE / credit limits); the current implementation
        uses risk-neutral for all outputs with a real-world scaling factor
        applied to PFE percentiles.
  \item Market factor dynamics follow the parametric models specified in
        Section~\ref{sec:theory} (Hull-White for rates, GBM for FX/equity, CIR
        for credit spreads).
  \item Netting sets are legally isolated: exposure netting is applied within
        but not across ISDA Master Agreements.
  \item The collateral posting party exercises its collateral obligations in
        full at each margin call date; no dispute or operational delay is
        modelled outside the Margin Period of Risk.
  \item SA-CCR calculation follows BCBS~279 (March 2014) as implemented in
        US rules effective April 2020.
\end{enumerate}

\subsection{Known Limitations}

\begin{enumerate}[noitemsep]
  \item The Gaussian copula used for joint factor simulation underestimates
        tail dependence in stress scenarios.
  \item The correlation matrix is static and recalibrated annually; intra-year
        correlation shifts are not captured.
  \item Gap risk for equity derivatives (overnight jumps) is not explicitly
        modelled.
  \item Approximately 8\% of counterparties require proxy CDS spread mapping,
        introducing basis risk.
  \item Path-dependent exotic derivatives use computational approximations that
        may introduce pricing errors of up to 5\%.
\end{enumerate}

% ─────────────────────────────────────────────────────────────────────────────
\section{Business Context and Purpose}
% ─────────────────────────────────────────────────────────────────────────────

\subsection{The Counterparty Credit Risk Problem}

OTC derivatives expose Apex to the risk that a counterparty defaults when the
derivative has positive mark-to-market value to Apex. Unlike a loan, where
exposure is the outstanding principal, a derivative's exposure is a random
variable that depends on the future evolution of market prices. The fundamental
challenge of CCR is that this exposure must be estimated over the entire life
of a contract --- potentially 30 years for a long-dated interest rate swap.

The 2008 financial crisis demonstrated the catastrophic consequences of
underestimating CCR. Bear Stearns and Lehman Brothers had large derivatives
portfolios with hundreds of counterparties. Exposure estimates that failed to
account for netting, or that used inadequate correlation assumptions, produced
dangerously low credit limits. When stress arrived, actual losses exceeded
modelled estimates by factors of 3--5x.

\subsection{Regulatory History}

CCR modelling has been shaped by successive waves of Basel regulation:

\begin{itemize}[noitemsep]
  \item \textbf{Basel II (2004)}: introduced the Internal Model Method (IMM)
        allowing banks with approved simulation models to use EPE for CCR
        capital, replacing the cruder Current Exposure Method (CEM).
  \item \textbf{Basel~III (2010)}: introduced the CVA capital charge requiring
        EPE-based CVA calculation; introduced the stressed EPE concept.
  \item \textbf{BCBS~279 / SA-CCR (2014)}: replaced CEM with the more
        risk-sensitive Standardised Approach to CCR, mandatory from 2020 for
        banks without IMM approval.
  \item \textbf{FRTB-CVA (2017)}: reformed the CVA capital framework,
        requiring banks to compute CVA sensitivities using an approved
        simulation infrastructure.
\end{itemize}

Apex currently uses SA-CCR for regulatory capital and maintains this model as
a candidate for IMM approval, which would allow use of internal EPE for
capital purposes.

\subsection{Model Users}

\begin{table}[h!]
\centering
\begin{tabular}{@{}lll@{}}
\toprule
\textbf{User} & \textbf{Function} & \textbf{Output Used} \\
\midrule
Credit Risk Officers   & Counterparty limit setting and monitoring & PFE at 97.5\% \\
XVA Desk               & CVA/FVA/MVA pricing and hedging           & EPE, ENE \\
Regulatory Capital     & Basel III SA-CCR RWA calculation          & SA-CCR EAD \\
Finance / P\&L         & CVA P\&L via downstream CVA model         & EPE \\
Internal Audit / MRO   & Model validation and ongoing monitoring   & All outputs \\
Regulators (OCC, Fed)  & CCAR stress testing; SR 11-7 review       & SA-CCR EAD, EPE \\
\bottomrule
\end{tabular}
\caption{Model users and primary outputs consumed}
\end{table}

% ─────────────────────────────────────────────────────────────────────────────
\section{Theoretical Foundation}
\label{sec:theory}
% ─────────────────────────────────────────────────────────────────────────────

\subsection{Close-Out Netting}

Under an ISDA Master Agreement, all transactions between two parties within a
single netting set are aggregated upon default. The close-out netting value at
time $t$ is:

\begin{equation}
  V_{\text{net}}(t) = \sum_{k=1}^{K} V_k(t)
  \label{eq:netting}
\end{equation}

where $V_k(t)$ is the mark-to-market value of transaction $k$ at time $t$,
measured from Apex's perspective (positive = Apex is owed money).

The \emph{netting benefit} at time $t$ quantifies how much the netting agreement
reduces gross positive exposure:

\begin{equation}
  \text{NettingBenefit}(t)
    = \sum_{k=1}^{K} \max(V_k(t),\, 0)
    - \max\!\left(\sum_{k=1}^{K} V_k(t),\, 0\right)
  \label{eq:netting_benefit}
\end{equation}

A netting benefit of zero implies either all trades are in the same direction
(no offset) or the netting set is entirely out-of-the-money for Apex.

\subsection{Collateral Offset and the CSA}

Where a Credit Support Annex (CSA) governs variation margin, the posted
collateral reduces the effective exposure at each simulation date. Let $C(t)$
denote the collateral balance (positive = collateral held by Apex), $\Theta$
the CSA threshold, and $\delta$ the minimum transfer amount (MTA). The
collateralised exposure before MPoR adjustment is:

\begin{equation}
  V_{\text{col}}(t)
    = \max\!\bigl(V_{\text{net}}(t) - C(t) - \Theta - \delta,\; 0\bigr)
  \label{eq:col_exposure}
\end{equation}

For a fully collateralised CSA with $\Theta = 0$ and $\delta = 0$, the exposure
collapses to $\max(V_{\text{net}}(t) - C(t), 0)$.

\subsection{Margin Period of Risk}

The Margin Period of Risk (MPoR) is the time from the last successful margin
call to the close-out of the portfolio following a counterparty default. During
MPoR, Apex cannot call additional margin; the last received collateral may be
stale. BCBS~261 specifies:

\begin{itemize}[noitemsep]
  \item \textbf{10 business days}: daily-margined netting sets (standard);
  \item \textbf{20 business days}: netting sets where the counterparty has
        $\geq 5{,}000$ transactions or illiquid collateral;
  \item \textbf{At least 20 business days}: re-hypothecated collateral.
\end{itemize}

The effective collateral at simulation time $t$, accounting for MPoR delay
$\tau$, is:

\begin{equation}
  C_{\text{eff}}(t) = C(t - \tau)
  \label{eq:coleff}
\end{equation}

where $\tau$ is the applicable MPoR. Substituting into \eqref{eq:col_exposure}:

\begin{equation}
  V_{\text{col,MPoR}}(t)
    = \max\!\bigl(V_{\text{net}}(t) - C(t-\tau) - \Theta - \delta,\; 0\bigr)
  \label{eq:mpor_exposure}
\end{equation}

\subsection{Exposure Metrics}

Let $\tilde{V}(t) \equiv V_{\text{col,MPoR}}(t)$ denote the collateralised,
MPoR-adjusted exposure. All exposure metrics are derived from the distribution
of $\tilde{V}(t)$ across Monte Carlo paths.

\subsubsection{Expected Exposure and Expected Positive / Negative Exposure}

\begin{align}
  \text{EE}(t)  &= \mathbb{E}\!\left[\max\!\bigl(\tilde{V}(t),\, 0\bigr)\right]
  \label{eq:ee} \\[4pt]
  \text{EPE}    &= \frac{1}{T} \int_0^T \text{EE}(t)\, dt
  \label{eq:epe} \\[4pt]
  \text{ENE}(t) &= \mathbb{E}\!\left[\min\!\bigl(\tilde{V}(t),\, 0\bigr)\right]
  \label{eq:ene}
\end{align}

The EPE in \eqref{eq:epe} represents the time-averaged expected exposure and
is the primary input to the CVA calculation.

\subsubsection{Effective EPE (Regulatory)}

For Basel IMM purposes, the \emph{Effective Expected Exposure} is the
non-decreasing envelope of EE, ensuring the regulatory capital measure is
conservative with respect to roll-off of near-term exposure:

\begin{equation}
  \text{EffEE}(t) = \max\!\bigl(\text{EE}(t),\; \text{EffEE}(t - \Delta t)\bigr),
  \quad \text{EffEE}(0) = V_{\text{net}}(0)
  \label{eq:eff_ee}
\end{equation}

The \emph{Effective EPE} is then:

\begin{equation}
  \text{EffEPE} = \frac{1}{1\text{Y}} \int_0^{1\text{Y}} \text{EffEE}(t)\, dt
  \label{eq:eff_epe}
\end{equation}

\subsubsection{Potential Future Exposure}

PFE at confidence level $\alpha$ at time horizon $t$ is the $\alpha$-quantile
of the exposure distribution:

\begin{equation}
  \text{PFE}_\alpha(t) = Q_\alpha\!\left[\tilde{V}(t)\right]
  \label{eq:pfe}
\end{equation}

Apex uses $\alpha = 97.5\%$ as the standard PFE confidence level for credit
limit purposes, consistent with Basel counterparty credit risk guidance.

\subsection{Market Factor Models}

\subsubsection{Interest Rates: Two-Factor Hull-White}

Domestic and foreign short rates are modelled using the two-factor Hull-White
(HW2) model, which fits the initial term structure exactly while allowing
mean-reverting stochastic evolution:

\begin{align}
  r(t) &= x_1(t) + x_2(t) + \varphi(t) \label{eq:hw2_r}\\
  dx_1(t) &= -a_1\, x_1(t)\, dt + \sigma_1\, dW_1(t) \label{eq:hw2_x1}\\
  dx_2(t) &= -a_2\, x_2(t)\, dt + \sigma_2\, dW_2(t) \label{eq:hw2_x2}
\end{align}

where $\varphi(t)$ is the deterministic fit to the initial forward curve,
$a_1, a_2$ are mean-reversion speeds, $\sigma_1, \sigma_2$ are volatilities,
and $\text{Corr}(dW_1, dW_2) = \rho_{12}$.

Current calibration (March 2026):
$a_1 = 0.03$, $a_2 = 0.35$, $\sigma_1 = 0.007$, $\sigma_2 = 0.014$,
$\rho_{12} = -0.72$.

\subsubsection{FX Rates: Geometric Brownian Motion}

FX spot rates $S_{\text{CCY/USD}}(t)$ are modelled as:

\begin{equation}
  \frac{dS(t)}{S(t)} = \bigl(r_d(t) - r_f(t)\bigr)\, dt + \sigma_{\text{FX}}\, dW_{\text{FX}}(t)
  \label{eq:fx_gbm}
\end{equation}

where $r_d$ and $r_f$ are domestic and foreign short rates from the HW2 models,
and $\sigma_{\text{FX}}$ is the FX volatility calibrated to ATM options with
maturity matching the simulation horizon.

\subsubsection{Equity: Heston Stochastic Volatility}

Equity prices are modelled under the Heston (1993) framework to capture the
volatility smile:

\begin{align}
  \frac{dS(t)}{S(t)} &= \mu\, dt + \sqrt{v(t)}\, dW_S(t) \label{eq:heston_s}\\
  dv(t) &= \kappa\bigl(\bar{v} - v(t)\bigr)\, dt + \xi\sqrt{v(t)}\, dW_v(t) \label{eq:heston_v}
\end{align}

with $\text{Corr}(dW_S, dW_v) = \rho_{\text{SV}}$. Parameters are calibrated
to the equity volatility surface.

\subsubsection{Credit Spreads: Cox-Ingersoll-Ross}

Credit spreads (used for CVA hazard rate evolution) follow a CIR process,
ensuring non-negativity:

\begin{equation}
  d\lambda(t) = \kappa_\lambda\!\bigl(\bar{\lambda} - \lambda(t)\bigr)\, dt
               + \sigma_\lambda\sqrt{\lambda(t)}\, dW_\lambda(t)
  \label{eq:cir}
\end{equation}

\subsection{Joint Simulation via Cholesky Decomposition}

All market factors are simulated jointly using a $47 \times 47$ correlation
matrix $\mathbf{P}$, Cholesky-decomposed into $\mathbf{L}$ such that
$\mathbf{L}\mathbf{L}^\top = \mathbf{P}$. Standard normal variates
$\mathbf{z} \sim \mathcal{N}(\mathbf{0}, \mathbf{I})$ are transformed into
correlated variates $\mathbf{w} = \mathbf{L}\mathbf{z}$, which are then
used to drive the stochastic differential equations above.

The 47 factors comprise: 5 interest rate curves (SOFR, ESTR, SONIA, TONAR,
HIBOR), 10 FX pairs (major and EM currencies), 12 equity indices, 15 credit
spread curves (investment grade, high yield, financials, sovereigns), and
5 commodity indices.

\subsection{SA-CCR: Standardised Approach to Counterparty Credit Risk}

For regulatory capital purposes under Basel~III (BCBS~279), the
Exposure-at-Default (EAD) for a netting set is:

\begin{equation}
  \text{EAD} = \alpha \cdot \bigl(\text{RC} + \text{PFE}_{\text{addon}}\bigr)
  \label{eq:sa_ccr_ead}
\end{equation}

where $\alpha = 1.4$ is the supervisory scalar, RC is the replacement cost, and
$\text{PFE}_{\text{addon}}$ is the supervisory add-on.

\subsubsection{Replacement Cost}

For margined netting sets:

\begin{equation}
  \text{RC} = \max\!\bigl(V - C,\; \text{TH} + \text{MTA} - \text{NICA},\; 0\bigr)
  \label{eq:rc_margined}
\end{equation}

For unmargined netting sets:

\begin{equation}
  \text{RC} = \max\!\bigl(V - C,\; 0\bigr)
  \label{eq:rc_unmargined}
\end{equation}

where $V$ is current net MtM, $C$ is collateral held net of haircuts, TH is
the CSA threshold, MTA is the minimum transfer amount, and NICA is the
net independent collateral amount.

\subsubsection{PFE Add-On Aggregation}

The aggregate PFE add-on across all five asset classes (interest rate, FX,
credit, equity, commodity) is:

\begin{equation}
  \text{PFE}_{\text{addon}} = \text{AddOn}^{\text{IR}}
    + \text{AddOn}^{\text{FX}}
    + \text{AddOn}^{\text{credit}}
    + \text{AddOn}^{\text{EQ}}
    + \text{AddOn}^{\text{com}}
  \label{eq:pfe_addon}
\end{equation}

Each asset-class add-on is computed from the supervisory delta, notional,
maturity factor, and supervisory volatility as specified in BCBS~279 Annex~B.
For interest rate derivatives, the add-on within a single currency hedging set
uses the partial netting formula:

\begin{equation}
  \text{AddOn}^{\text{IR, CCY}}
    = \left|\, \sum_j \delta_j \cdot d_j \cdot \text{MF}_j \cdot \sigma^{\text{IR}}_j
      \right|
  \label{eq:ir_addon}
\end{equation}

where $\delta_j \in \{+1, -1\}$ is the supervisory delta, $d_j$ is the
adjusted notional, $\text{MF}_j$ is the maturity factor
($\text{MF}_j = \sqrt{\min(M_j, 1\text{Y}) / 1\text{Y}}$ for unmargined), and
$\sigma^{\text{IR}}_j$ is the supervisory volatility ($= 0.50\%$ for IR).

% ─────────────────────────────────────────────────────────────────────────────
\section{Data Requirements and Governance}
% ─────────────────────────────────────────────────────────────────────────────

\subsection{Input Data Sources}

\begin{longtable}{@{}lll@{}}
\toprule
\textbf{Data Item} & \textbf{Source} & \textbf{Frequency} \\
\midrule
\endhead
\bottomrule
\endfoot
OIS/RFR yield curves (SOFR, ESTR, SONIA, TONAR, HIBOR) & Bloomberg / Internal & Daily \\
FX spot rates (10 pairs vs. USD) & Bloomberg & Real-time (EOD snap) \\
Equity index levels (12 indices) & Bloomberg & Daily \\
Credit spread curves (IG, HY, fins, sov) & Markit / Bloomberg & Daily \\
Swaption vol surfaces (all rate curves) & Bloomberg & Daily \\
FX ATM vol (all pairs) & Bloomberg & Daily \\
Equity vol surfaces (all indices) & Bloomberg & Daily \\
Credit spread vol (CDS options where available) & Markit & Weekly \\
ISDA Master Agreement database & Legal / ISDA operations & On amendment \\
CSA terms (threshold, MTA, eligible collateral) & Legal / ISDA operations & On amendment \\
Trade population (all OTC derivatives) & Front-office systems & Real-time \\
Collateral balances (VM posted/received) & Collateral Management & Daily \\
Historical factor data (5-year lookback) & Bloomberg / Proprietary & Annual \\
\end{longtable}

\subsection{Data Quality Controls}

\begin{enumerate}[noitemsep]
  \item \textbf{Curve completeness}: All OIS curves must have at least 10 tenor
        points covering 1W to 30Y. Missing tenors use linear log-discount
        interpolation with a staleness flag if the bracketing points are
        $>2$ business days old.
  \item \textbf{Vol surface consistency}: Swaption vol surfaces are validated
        for absence of calendar spread arbitrage and butterfly spread arbitrage
        before use. Surfaces failing validation revert to prior-day with an
        alert.
  \item \textbf{Trade data completeness}: Trade population must reconcile to
        front-office systems within 0.5\% of notional before the simulation
        run proceeds. Material breaks ($>1\%$) halt the run and escalate to
        Operations.
  \item \textbf{CSA database}: The CSA terms database must cover $\geq 99\%$
        of live netting sets by notional. Missing CSA terms default to
        ``uncollateralised'' (conservative).
\end{enumerate}

\subsection{BCBS 239 Compliance}

As a Tier~1 model feeding regulatory capital, the PFE/CCR engine is subject to
BCBS~239 data lineage requirements. Every input data item is tagged with
provenance metadata (source system, timestamp, transformation applied).
The full data lineage from raw market data to SA-CCR EAD must be
reconstructable within 4 hours of a regulatory request.

% ─────────────────────────────────────────────────────────────────────────────
\section{Methodology and Implementation}
% ─────────────────────────────────────────────────────────────────────────────

\subsection{Monte Carlo Engine}

The simulation engine generates $N = 10{,}000$ risk-neutral paths over
$M = 200$ time steps:

\begin{itemize}[noitemsep]
  \item $t_1, \ldots, t_{52}$: weekly steps to 1 year;
  \item $t_{53}, \ldots, t_{112}$: monthly steps from 1 to 5 years;
  \item $t_{113}, \ldots, t_{200}$: quarterly steps from 5 to 30 years.
\end{itemize}

GPU acceleration (CUDA via cuBLAS) is used for the $47 \times 47$ Cholesky
decomposition and for batched path generation. The full portfolio run completes
in under 4 hours on the production GPU cluster. An intraday ``Greeks'' run
(500 paths, 1Y horizon only) completes in under 30 minutes.

\subsection{Netting Set Aggregation Algorithm}

For each simulation path $\omega$ and time step $t$:

\begin{enumerate}[noitemsep]
  \item Price all trades $V_k(\omega, t)$ using the simulated market factors.
  \item Group trades by ISDA netting set identifier.
  \item Compute $V_{\text{net}}(\omega, t) = \sum_k V_k(\omega, t)$ within
        each netting set (Equation~\ref{eq:netting}).
  \item Apply collateral offset per Equation~\ref{eq:mpor_exposure}, using
        the CSA terms for the netting set.
  \item Record $\tilde{V}(\omega, t)$ as the path-level collateralised exposure.
\end{enumerate}

\subsection{Exposure Profile Construction}

After all $N$ paths are generated, for each netting set and time step $t$:

\begin{enumerate}[noitemsep]
  \item $\text{EE}(t) = \frac{1}{N} \sum_\omega \max(\tilde{V}(\omega, t), 0)$;
  \item $\text{ENE}(t) = \frac{1}{N} \sum_\omega \min(\tilde{V}(\omega, t), 0)$;
  \item $\text{PFE}_{97.5\%}(t) = $ 97.5th-percentile of
        $\{\tilde{V}(\omega, t)\}_{\omega=1}^N$;
  \item $\text{EffEE}(t)$ computed recursively per Equation~\ref{eq:eff_ee}.
\end{enumerate}

Portfolio-level profiles aggregate netting-set profiles using the same
path-consistent summation (no diversification benefit across netting sets for
EPE, consistent with Basel IMM requirements).

\subsection{SA-CCR Calculation}

The SA-CCR module runs as a post-processing step on the current-date trade
population (no simulation paths required). For each netting set:

\begin{enumerate}[noitemsep]
  \item Classify each trade into asset class (IR, FX, credit, EQ, commodity).
  \item Compute supervisory delta $\delta_j$ for each trade (sign and optionality
        adjustment).
  \item Compute adjusted notional $d_j$ (trade-specific; tenor-adjusted for IR).
  \item Compute maturity factor $\text{MF}_j$ (margined or unmargined formula).
  \item Aggregate add-ons within hedging sets (with partial netting for IR/credit).
  \item Aggregate hedging-set add-ons to netting-set $\text{PFE}_{\text{addon}}$.
  \item Compute RC per Equation~\ref{eq:rc_margined} or \ref{eq:rc_unmargined}.
  \item Compute EAD per Equation~\ref{eq:sa_ccr_ead}.
\end{enumerate}

\subsection{Shared Infrastructure with CVA, FVA, MVA}

CVA, FVA, and MVA all consume outputs from the same Monte Carlo run. This
ensures consistency across XVA measures --- a requirement under IFRS~13 for
coherent fair value measurement --- and reduces computational cost by avoiding
redundant simulation. The dependency chain is:

\begin{center}
  \textbf{PFE/CCR Engine} $\;\longrightarrow\;$ \textbf{CVA} (APEX-MDL-0010)\\[2pt]
  \textbf{PFE/CCR Engine} $\;\longrightarrow\;$ \textbf{FVA} (APEX-MDL-0011)\\[2pt]
  \textbf{PFE/CCR Engine} $\;\longrightarrow\;$ \textbf{MVA} (APEX-MDL-0012)
\end{center}

A failure of the PFE/CCR engine propagates to all three downstream models.
The fallback procedure (prior-day EPE/ENE scaled by Greeks) is documented
in Section~\ref{sec:controls}.

% ─────────────────────────────────────────────────────────────────────────────
\section{Model Testing and Backtesting}
% ─────────────────────────────────────────────────────────────────────────────

\subsection{P\&L Explain for EE}

Daily predicted EE change is compared to actual EE change:

\begin{equation}
  \Delta\text{EE}_{\text{predicted}}
    = \frac{\partial \text{EE}}{\partial r} \Delta r
    + \frac{\partial \text{EE}}{\partial S_{\text{FX}}} \Delta S_{\text{FX}}
    + \frac{\partial \text{EE}}{\partial \lambda} \Delta\lambda
    + \text{theta}
\end{equation}

P\&L explain ratio $= \Delta\text{EE}_{\text{predicted}} / \Delta\text{EE}_{\text{actual}}$
must exceed 85\%. Deviations $>15\%$ trigger a model investigation within
2 business days.

\subsection{SA-CCR Benchmark}

The SA-CCR EAD is compared quarterly to the 97.5th-percentile internal PFE
multiplied by $\alpha = 1.4$ (the SA-CCR scalar):

\begin{equation}
  \text{SA-CCR ratio}
    = \frac{\text{SA-CCR EAD}}{1.4 \times \text{PFE}_{97.5\%}(1\text{Y})}
\end{equation}

Acceptable range: $[0.9,\; 1.5]$. Values outside this range indicate either
over-conservatism in SA-CCR or inadequacy of the internal model, and trigger
a 30-day review.

\subsection{Historical Backtesting}

EE predictions are compared to realised exposures on a rolling 1-year
lookback. The test asks: ``did the realised 1-year maximum exposure fall
below the predicted PFE$_{97.5\%}$ at least 97.5\% of the time?'' Following
the Basel traffic light framework:

\begin{table}[h!]
\centering
\begin{tabular}{@{}lll@{}}
\toprule
\textbf{Zone} & \textbf{Exceptions (out of 250 obs.)} & \textbf{Action} \\
\midrule
Green  & 0--6   & No action required \\
Yellow & 7--9   & Explanation required; potential add-on \\
Red    & 10+    & Model review mandatory; capital add-on \\
\bottomrule
\end{tabular}
\caption{Backtesting traffic light zones (adapted for CCR)}
\end{table}

\subsection{Stress Scenario Analysis}

Five stress scenarios are applied to the current portfolio and the resulting
EE profiles are reported to the Market Risk Officer quarterly:

\begin{table}[h!]
\centering
\begin{tabular}{@{}lrr@{}}
\toprule
\textbf{Scenario} & \textbf{EE Increase (\$M)} & \textbf{PFE$_{97.5\%}$ Increase (\$M)} \\
\midrule
2008 GFC replay              & +\$1{,}100 & +\$2{,}400 \\
2020 COVID stress            & +\$650     & +\$1{,}500 \\
Rates +200bp parallel shift  & +\$420     & +\$980 \\
Credit spreads +100bp        & +\$280     & +\$620 \\
FX +/-20\% (dollar weakening) & +\$210     & +\$450 \\
\bottomrule
\end{tabular}
\caption{Stress scenario EE and PFE impacts (approximate, current portfolio)}
\end{table}

% ─────────────────────────────────────────────────────────────────────────────
\section{Performance Metrics and Calibration Standards}
% ─────────────────────────────────────────────────────────────────────────────

\begin{longtable}{@{}llll@{}}
\toprule
\textbf{Metric} & \textbf{Acceptable Range} & \textbf{Frequency} & \textbf{Action if Breached} \\
\midrule
\endhead
\bottomrule
\endfoot
EE P\&L Explain                & $>85\%$           & Daily    & Investigate within 2 business days \\
SA-CCR ratio (SA-CCR / 1.4$\times$PFE) & $[0.9,\, 1.5]$  & Quarterly & Review within 30 days \\
Backtesting exceptions (annual) & $\leq 6$ (Green zone) & Quarterly & Escalate if Yellow; mandatory if Red \\
Correlation matrix condition \# & $<1{,}000$        & Annual   & Recalibrate; notify MRO \\
Data completeness (CSA coverage) & $\geq 99\%$       & Daily    & Escalate to Legal/Operations \\
Intraday run completion time    & $<30$ minutes      & Daily    & Technology escalation \\
Full EOD run completion time    & $<4$ hours         & Daily    & Technology escalation \\
\end{longtable}

\subsection{Recalibration Schedule}

\begin{itemize}[noitemsep]
  \item \textbf{HW2 and CIR parameters}: Annual recalibration to 5-year
        historical data; ad hoc recalibration upon material market structure
        change (e.g., rate regime change).
  \item \textbf{Heston equity parameters}: Quarterly recalibration to current
        vol surface.
  \item \textbf{Correlation matrix}: Annual recalibration using rolling 5-year
        equally-weighted data; reviewed after market stress events.
  \item \textbf{SA-CCR supervisory parameters}: Fixed by BCBS~279; updated
        only upon regulatory change with full MRO review.
\end{itemize}

% ─────────────────────────────────────────────────────────────────────────────
\section{Limitations, Assumptions, and Known Weaknesses}
% ─────────────────────────────────────────────────────────────────────────────

\subsection{Gaussian Copula Tail Dependence}

The joint factor simulation uses a Gaussian copula: correlated normal random
variables are mapped to each factor's marginal distribution. The Gaussian
copula underestimates tail dependence --- the probability of multiple factors
simultaneously moving to extreme values is lower in the Gaussian copula than
in empirical data, particularly during systemic stress events.

\textbf{Compensating control}: Stressed EE (computed using 2008 and 2020 shock
scenarios) is disclosed alongside base EE. A capital add-on for tail dependence
is included in the model reserve.

\subsection{Static Correlation Matrix}

The 47$\times$47 correlation matrix is calibrated once per year. During periods
of market stress, correlations can change rapidly and materially (e.g.,
equity-FX correlations invert during risk-off events). The static matrix may
produce inaccurate exposure estimates intra-year.

\textbf{Compensating control}: The Market Risk Officer monitors intraday
correlation estimates from a 60-day rolling window and can trigger an
unscheduled recalibration if any major pair moves $>30$ percentage points.

\subsection{Equity Gap Risk}

The Heston model uses continuous-time diffusion with no jump component.
Overnight equity gaps (e.g., a 20\% single-day drop following an unexpected
corporate event) are not captured. For equity derivatives with short-dated
optionality, this can understate PFE.

\textbf{Compensating control}: Equity derivatives with remaining maturity
$<6$ months use a jump-adjusted vol ($\sigma_{\text{adj}} = 1.2 \times
\sigma_{\text{Heston}}$) for PFE calculation only.

\subsection{Proxy Credit Spread Mapping}

Approximately 8\% of counterparties by notional have no observable CDS spread.
These use a proxy mapping based on credit rating and sector (e.g., BBB-rated
European utility $\mapsto$ iTraxx Main component average). The proxy spread
may not reflect idiosyncratic credit risk.

\textbf{Compensating control}: Proxy-mapped counterparties are flagged in all
outputs. A 25\% add-on is applied to the EPE for proxy-mapped counterparties
in the CVA calculation. The Credit Risk function reviews proxy assignments
quarterly.

\subsection{Exotic Derivative Approximations}

Path-dependent exotic derivatives (e.g., barrier options, Asian options, TARFs)
are priced using closed-form approximations or pre-computed exposure grids
rather than full nested Monte Carlo. These approximations may introduce pricing
errors of up to 5\% for individual instruments.

\textbf{Compensating control}: Exotics exceeding \$50M in notional receive a
full nested Monte Carlo re-price on a monthly basis. The largest 20 exotic
positions are benchmarked to an independent pricing library quarterly.

% ─────────────────────────────────────────────────────────────────────────────
\section{Compensating Controls}
\label{sec:controls}
% ─────────────────────────────────────────────────────────────────────────────

\begin{longtable}{@{}lll@{}}
\toprule
\textbf{Risk} & \textbf{Control} & \textbf{Owner} \\
\midrule
\endhead
\bottomrule
\endfoot
Tail dependence underestimation
  & Stressed EE disclosure; model reserve add-on
  & XVA Quant / MRO \\
Static correlation
  & 60-day rolling monitor; unscheduled recalibration trigger at 30pp move
  & Market Risk \\
Equity gap risk
  & 1.2$\times$ vol adjustment for short-dated equity derivatives
  & XVA Quant \\
Proxy CDS mapping
  & 25\% EPE add-on; quarterly review by Credit Risk
  & Credit Risk \\
Exotic approximation error
  & Monthly full reprice for $>$\$50M; quarterly benchmark vs.\ independent library
  & XVA Technology \\
Engine failure (downstream cascade)
  & Fallback: prior-day EPE/ENE scaled by Greeks; alert to CVA/FVA/MVA desks
  & XVA Technology \\
CSA database gaps
  & Missing CSA defaults to uncollateralised (conservative); nightly reconciliation
  & Legal / Operations \\
Data staleness
  & 2-business-day staleness alert; fallback to prior-day with flag
  & Market Data \\
\end{longtable}

\subsection{Model Reserve}

A model reserve of \textbf{\$35 million} is held against PFE/CCR model
uncertainty:

\begin{itemize}[noitemsep]
  \item Gaussian copula tail dependence: \$15M;
  \item Proxy CDS mapping basis risk: \$10M;
  \item Exotic derivative path-dependency approximation: \$10M.
\end{itemize}

The reserve is reviewed annually alongside the model recalibration cycle and
may be adjusted following material market events.

% ─────────────────────────────────────────────────────────────────────────────
\section{Relationship to Downstream Models}
% ─────────────────────────────────────────────────────────────────────────────

The PFE/CCR Exposure Engine is a producer model; it does not consume outputs
from other models within the XVA suite. The table below summarises the
downstream dependencies:

\begin{table}[h!]
\centering
\begin{tabular}{@{}llll@{}}
\toprule
\textbf{Downstream Model} & \textbf{Model ID} & \textbf{Input Consumed} & \textbf{Sensitivity} \\
\midrule
CVA  & APEX-MDL-0010 & EPE profile (all counterparties)       & High --- linear in EPE \\
FVA  & APEX-MDL-0011 & EPE and ENE profiles                   & High --- linear in EPE/ENE \\
MVA  & APEX-MDL-0012 & Simulated IM profiles (Monte Carlo)    & High --- direct path dependency \\
ColVA & APEX-MDL-0013 & $|\text{EPE}|$ for collateral balance & Medium \\
\bottomrule
\end{tabular}
\caption{Downstream model dependencies}
\end{table}

Any change to the PFE/CCR engine that affects EPE or ENE profiles (e.g.,
correlation matrix recalibration, factor model change, netting set
re-aggregation) triggers a mandatory impact assessment for all four downstream
models before the change is deployed to production.

% ─────────────────────────────────────────────────────────────────────────────
\section{Change Management}
% ─────────────────────────────────────────────────────────────────────────────

\textbf{Material changes} requiring full MRO re-validation:

\begin{itemize}[noitemsep]
  \item Change to any market factor model (HW2 $\to$ LMM, Heston $\to$ SABR, etc.);
  \item Change to netting set aggregation logic or MPoR treatment;
  \item Change to SA-CCR add-on calculation methodology;
  \item Integration of new asset class or new market (e.g., adding crypto
        derivative exposure);
  \item Migration to a new simulation infrastructure or hardware platform.
\end{itemize}

\textbf{Non-material changes} via expedited review (5 business days):

\begin{itemize}[noitemsep]
  \item Annual recalibration of factor model parameters (within pre-approved bounds);
  \item Addition of a new counterparty to the proxy CDS mapping table;
  \item Computational optimisation with no change to outputs;
  \item Bug fixes that correct errors without changing methodology.
\end{itemize}

% ─────────────────────────────────────────────────────────────────────────────
\section{Approvals and Sign-Off}
% ─────────────────────────────────────────────────────────────────────────────

\begin{table}[h!]
\centering
\begin{tabular}{@{}llll@{}}
\toprule
\textbf{Role} & \textbf{Name} & \textbf{Status} & \textbf{Date} \\
\midrule
Model Developer      & XVA Quant Team      & $\checkmark$ Submitted & March 2026 \\
Model Risk Officer   & Dr.\ Rebecca Chen   & $\circ$ Pending & --- \\
Chief Risk Officer   & Dr.\ Priya Nair     & $\circ$ Pending & --- \\
Business Owner       & James Okafor        & $\circ$ Pending & --- \\
\bottomrule
\end{tabular}
\caption{Approval status}
\end{table}

\vspace{2em}
\hrule
\vspace{0.5em}
{\small\color{apexgray}
Document Classification: RESTRICTED --- Model Risk Sensitive\\
Apex Global Bank \textbar\ Quantitative Research Division \textbar\ XVA Team\\
Next scheduled review: March 2027 (or upon material model change)
}

% ─────────────────────────────────────────────────────────────────────────────
\appendix
\section{Notation Summary}
% ─────────────────────────────────────────────────────────────────────────────

\begin{longtable}{@{}ll@{}}
\toprule
\textbf{Symbol} & \textbf{Definition} \\
\midrule
\endhead
\bottomrule
\endfoot
$V_k(t)$          & MtM value of trade $k$ at time $t$ (Apex's perspective) \\
$V_{\text{net}}(t)$ & Net MtM value of netting set at time $t$ \\
$C(t)$            & Collateral balance at time $t$ (positive = held by Apex) \\
$\Theta$          & CSA threshold \\
$\delta$          & CSA minimum transfer amount (MTA) \\
$\tau$            & Margin Period of Risk \\
$\tilde{V}(t)$    & Collateralised, MPoR-adjusted exposure \\
$\text{EE}(t)$    & Expected Exposure at time $t$ \\
$\text{EPE}$      & (Time-averaged) Expected Positive Exposure \\
$\text{EffEE}(t)$ & Effective Expected Exposure (non-decreasing) \\
$\text{ENE}(t)$   & Expected Negative Exposure at time $t$ \\
$\text{PFE}_\alpha(t)$ & Potential Future Exposure at confidence $\alpha$, horizon $t$ \\
$\alpha$          & SA-CCR supervisory scalar ($= 1.4$) \\
$\text{RC}$       & SA-CCR Replacement Cost \\
$\text{EAD}$      & Exposure at Default (SA-CCR) \\
$N$               & Number of Monte Carlo paths ($= 10{,}000$) \\
$M$               & Number of simulation time steps ($= 200$) \\
$a_1, a_2$        & Hull-White mean-reversion speeds \\
$\sigma_1, \sigma_2$ & Hull-White volatilities \\
$\kappa, \bar{v}, \xi$ & Heston mean-reversion, long-run variance, vol-of-vol \\
$\kappa_\lambda, \bar{\lambda}, \sigma_\lambda$ & CIR credit spread parameters \\
$\mathbf{P}$      & $47 \times 47$ factor correlation matrix \\
$\mathbf{L}$      & Cholesky factor of $\mathbf{P}$ \\
\end{longtable}

% ─────────────────────────────────────────────────────────────────────────────
\end{document}
